\documentclass[12pt,a4paper]{article}

% ===== PACKAGES =====
\usepackage{amsmath}
\usepackage{amssymb}
\usepackage{geometry}
\geometry{margin=2.5cm}
\usepackage{booktabs}
\usepackage{xcolor}
\usepackage{hyperref}
\usepackage{listings}
\usepackage{pgfplots}
\pgfplotsset{compat=1.18}
\usepackage[most,breakable]{tcolorbox}
\tcbuselibrary{skins, breakable}
\usepackage{enumitem}
\usepackage{fancyhdr}
\usepackage{tikz}
\usepackage{array}

% ===== TCOLORBOX ENVIRONMENTS =====
\newtcolorbox{curiositybox}[1]{colback=yellow!10, colframe=orange!80!black, fonttitle=\bfseries, title=#1, breakable}
\newtcolorbox{infobox}[1]{colback=blue!5, colframe=blue!60!black, fonttitle=\bfseries, title=#1, breakable}
\newtcolorbox{mistakebox}[1]{colback=red!5, colframe=red!60!black, fonttitle=\bfseries, title=#1, breakable}
\newtcolorbox{learnbox}[1]{colback=green!5, colframe=green!60!black, fonttitle=\bfseries, title=#1, breakable}

% ===== LSTLISTING SETTINGS =====
\lstset{language=Python, basicstyle=\ttfamily\small, keywordstyle=\color{blue},
        commentstyle=\color{gray}, stringstyle=\color{red},
        showstringspaces=false, breaklines=true, frame=single}

% ===== FANCYHDR =====
\pagestyle{fancy}
\fancyhf{}
\lhead{Topic 08: Milne Thomson Method}
\rhead{Unit I --- Complex Variables}
\cfoot{\thepage}

% ===== TITLE =====
\title{\textbf{Topic 08: Construction of Analytic Function\\by Milne Thomson Method}\\[6pt]
\large Unit I --- Complex Variable: Differentiation\\
Complex Variables and Special Functions\\
JNTUA College of Engineering (Autonomous)}
\author{}
\date{}

% ===== DOCUMENT =====
\begin{document}
\maketitle
\tableofcontents
\newpage

% ------------------------------------------------------------------
% OPENING HOOK
% ------------------------------------------------------------------
\begin{curiositybox}{Hook}
Up to now, finding an analytic function often meant first hunting down the harmonic conjugate. Milne Thomson's method offers a shortcut: instead of reconstructing the missing part step by step, it builds the analytic function directly.
\end{curiositybox}

\vspace{0.4cm}

% ------------------------------------------------------------------
% SECTION 1
% ------------------------------------------------------------------
\section{Why This Method Is Useful}

In earlier work, constructing an analytic function \( f(z) = u + iv \) from a known harmonic part required integrating partial derivatives to find the conjugate, then combining. While systematic, that approach accumulates integration constants and demands careful bookkeeping at every step.

Milne Thomson's method bypasses the intermediate harmonic-conjugate construction entirely. Once the real part \( u(x,y) \) or imaginary part \( v(x,y) \) is known, one algebraic substitution --- replacing \( x \) with \( z \) and \( y \) with \( 0 \) --- yields \( f(z) \) directly in terms of \( z \) alone, with no partial integrations required.

\begin{center}
\begin{tabular}{@{}lll@{}}
\toprule
\textbf{Step} & \textbf{Harmonic Conjugate Method} & \textbf{Milne Thomson Method} \\
\midrule
Start & Know \( u(x,y) \) & Know \( u(x,y) \) \\
Core operation & Integrate \( \partial u/\partial x \), \( \partial u/\partial y \) & Compute \( \phi_1 \), \( \phi_2 \) then substitute \\
Integration & Required (may be involved) & Not required \\
Final form & \( f(z) = u + iv \) assembled & \( f(z) \) written directly in \( z \) \\
Pitfall & Integration constant bookkeeping & Must verify harmonicity first \\
\bottomrule
\end{tabular}
\end{center}

\begin{learnbox}{Key Insight}
The power of Milne Thomson's method lies in converting a two-variable reconstruction into a single-variable expression. It does not remove the need to check that \( u \) (or \( v \)) is harmonic before applying the method.
\end{learnbox}

% ------------------------------------------------------------------
% SECTION 2
% ------------------------------------------------------------------
\section{Background Needed}

\begin{infobox}{Essential Prerequisites}
\textbf{Analytic function:} \( f(z) = u(x,y) + iv(x,y) \) is analytic in a domain \( D \) if it is differentiable at every point of \( D \).

\textbf{Cauchy-Riemann equations (C-R):}
\[
\frac{\partial u}{\partial x} = \frac{\partial v}{\partial y}, \qquad \frac{\partial u}{\partial y} = -\frac{\partial v}{\partial x}.
\]

\textbf{Harmonic function:} A real-valued function \( \phi(x,y) \) satisfying Laplace's equation
\[
\nabla^2 \phi = \frac{\partial^2 \phi}{\partial x^2} + \frac{\partial^2 \phi}{\partial y^2} = 0.
\]

\textbf{Harmonic conjugate:} If \( u \) is harmonic, there exists \( v \) (unique up to a constant) such that \( f = u + iv \) is analytic; \( v \) is the harmonic conjugate of \( u \).

\textbf{Complex derivative:}
\[
f'(z) = \frac{\partial u}{\partial x} + i\frac{\partial v}{\partial x} = \frac{\partial v}{\partial y} - i\frac{\partial u}{\partial y}.
\]
\end{infobox}

\begin{learnbox}{What to Verify First}
Before applying any construction method, \textbf{always} check that the given function satisfies Laplace's equation. If it does not, no analytic function with that real or imaginary part exists.
\end{learnbox}

% ------------------------------------------------------------------
% SECTION 3
% ------------------------------------------------------------------
\section{Milne Thomson Method When \texorpdfstring{\( u(x,y) \)}{u(x,y)} Is Given}

\begin{infobox}{Derivation --- Given \( u \)}
Let \( f(z) = u(x,y) + iv(x,y) \) be analytic. Then:
\[
f'(z) = \frac{\partial u}{\partial x} - i\frac{\partial u}{\partial y}.
\]
Define two real functions:
\[
\phi_1(x,y) = \frac{\partial u}{\partial x}, \qquad \phi_2(x,y) = \frac{\partial u}{\partial y}.
\]
So:
\[
f'(z) = \phi_1(x,y) - i\,\phi_2(x,y).
\]
\textbf{Key substitution:} Since \( z = x + iy \), setting \( y = 0 \) gives \( x = z \). Replace \( x \) by \( z \) and \( y \) by \( 0 \):
\[
f'(z) = \phi_1(z,0) - i\,\phi_2(z,0).
\]
Integrate with respect to \( z \):
\[
\boxed{f(z) = \int \bigl[\phi_1(z,0) - i\,\phi_2(z,0)\bigr]\,dz + C,}
\]
where \( C \) is an arbitrary complex constant.

\textbf{Why \( x = z, y = 0 \) works:}
The expression \( f'(z) = \phi_1(x,y) - i\phi_2(x,y) \) holds for all \( z = x+iy \). In particular it holds along the real axis (\( y=0 \), \( x = z \)), and by analytic continuation the formula \( f'(z) = \phi_1(z,0) - i\phi_2(z,0) \) extends to the whole domain.
\end{infobox}

\begin{learnbox}{Formula to Remember --- Given \( u \)}
\[
f(z) = \int \left[\frac{\partial u}{\partial x}\bigg|_{x=z,\,y=0} - i\,\frac{\partial u}{\partial y}\bigg|_{x=z,\,y=0}\right]dz + C.
\]
\end{learnbox}

% ------------------------------------------------------------------
% SECTION 4
% ------------------------------------------------------------------
\section{Milne Thomson Method When \texorpdfstring{\( v(x,y) \)}{v(x,y)} Is Given}

\begin{infobox}{Derivation --- Given \( v \)}
If \( v(x,y) \) is given, use:
\[
f'(z) = \frac{\partial v}{\partial y} + i\frac{\partial v}{\partial x}.
\]
Define:
\[
\psi_1(x,y) = \frac{\partial v}{\partial y}, \qquad \psi_2(x,y) = \frac{\partial v}{\partial x}.
\]
So:
\[
f'(z) = \psi_1(x,y) + i\,\psi_2(x,y).
\]
Substitute \( x = z,\; y = 0 \):
\[
\boxed{f(z) = \int \bigl[\psi_1(z,0) + i\,\psi_2(z,0)\bigr]\,dz + C.}
\]
\end{infobox}

\begin{learnbox}{Formula to Remember --- Given \( v \)}
\[
f(z) = \int \left[\frac{\partial v}{\partial y}\bigg|_{x=z,\,y=0} + i\,\frac{\partial v}{\partial x}\bigg|_{x=z,\,y=0}\right]dz + C.
\]
Note the sign pattern: the \( v \)-case uses \( +i \), the \( u \)-case uses \( -i \).
\end{learnbox}

% ------------------------------------------------------------------
% SECTION 5
% ------------------------------------------------------------------
\section{Step-by-Step Algorithm}

\begin{infobox}{Algorithm: Milne Thomson Method}
\textbf{Case A --- Given \( u(x,y) \):}
\begin{enumerate}[label=\textbf{Step \arabic*.}]
  \item Verify \( \nabla^2 u = 0 \) (Laplace's equation). If not satisfied, stop.
  \item Compute \( \phi_1(x,y) = \dfrac{\partial u}{\partial x} \) and \( \phi_2(x,y) = \dfrac{\partial u}{\partial y} \).
  \item Substitute \( x \to z \) and \( y \to 0 \) in both expressions.
  \item Form \( f'(z) = \phi_1(z,0) - i\,\phi_2(z,0) \).
  \item Integrate: \( f(z) = \displaystyle\int f'(z)\,dz + C \).
  \item Verify by expanding \( f(z) \) and checking that its real part equals \( u(x,y) \).
\end{enumerate}

\vspace{4pt}
\textbf{Case B --- Given \( v(x,y) \):}
\begin{enumerate}[label=\textbf{Step \arabic*.}]
  \item Verify \( \nabla^2 v = 0 \).
  \item Compute \( \psi_1(x,y) = \dfrac{\partial v}{\partial y} \) and \( \psi_2(x,y) = \dfrac{\partial v}{\partial x} \).
  \item Substitute \( x \to z \) and \( y \to 0 \).
  \item Form \( f'(z) = \psi_1(z,0) + i\,\psi_2(z,0) \).
  \item Integrate: \( f(z) = \displaystyle\int f'(z)\,dz + C \).
  \item Verify that the imaginary part of \( f(z) \) equals \( v(x,y) \).
\end{enumerate}
\end{infobox}

\begin{learnbox}{Memory Aid}
\textbf{u-given:} \( f'(z) = u_x - i\,u_y \) then substitute \( (z,0) \).\\
\textbf{v-given:} \( f'(z) = v_y + i\,v_x \) then substitute \( (z,0) \).
\end{learnbox}

% ------------------------------------------------------------------
% SECTION 6: WORKED EXAMPLES
% ------------------------------------------------------------------
\section{Worked Examples}

% ---- Example 1 ----
\subsection*{Example 1: Given \( u = x^2 - y^2 \)}

\textbf{Step 1 --- Verify harmonicity:}
\[
\frac{\partial^2 u}{\partial x^2} = 2, \qquad \frac{\partial^2 u}{\partial y^2} = -2, \qquad \nabla^2 u = 2 + (-2) = 0. \quad \checkmark
\]
\textbf{Step 2 --- Compute partials:}
\[
\phi_1 = \frac{\partial u}{\partial x} = 2x, \qquad \phi_2 = \frac{\partial u}{\partial y} = -2y.
\]
\textbf{Step 3 --- Substitute \( x=z,\; y=0 \):}
\[
\phi_1(z,0) = 2z, \qquad \phi_2(z,0) = 0.
\]
\textbf{Step 4 --- Form \( f'(z) \):}
\[
f'(z) = 2z - i(0) = 2z.
\]
\textbf{Step 5 --- Integrate:}
\[
f(z) = \int 2z\,dz + C = z^2 + C.
\]
\textbf{Verification:} \( f(z) = z^2 = (x+iy)^2 = x^2 - y^2 + 2ixy \). Real part \( = x^2 - y^2 = u \). \checkmark

\begin{learnbox}{Result}
\( f(z) = z^2 + C \). The real part \( x^2 - y^2 \) corresponds to the simplest non-trivial analytic function.
\end{learnbox}

% ---- Example 2 ----
\subsection*{Example 2: Given \( u = e^x \cos y \)}

\textbf{Step 1 --- Verify harmonicity:}
\[
\frac{\partial^2 u}{\partial x^2} = e^x \cos y, \qquad \frac{\partial^2 u}{\partial y^2} = -e^x \cos y, \qquad \nabla^2 u = 0. \quad \checkmark
\]
\textbf{Step 2 --- Compute partials:}
\[
\phi_1 = e^x \cos y, \qquad \phi_2 = -e^x \sin y.
\]
\textbf{Step 3 --- Substitute \( x=z,\; y=0 \):}
\[
\phi_1(z,0) = e^z \cos 0 = e^z, \qquad \phi_2(z,0) = -e^z \sin 0 = 0.
\]
\textbf{Step 4:}
\[
f'(z) = e^z - i(0) = e^z.
\]
\textbf{Step 5:}
\[
f(z) = \int e^z\,dz + C = e^z + C.
\]
\textbf{Verification:} \( e^z = e^{x+iy} = e^x(\cos y + i\sin y) \). Real part \( = e^x\cos y = u \). \checkmark

\begin{learnbox}{Result}
\( f(z) = e^z + C \). The exponential function is its own complex derivative.
\end{learnbox}

% ---- Example 3 ----
\subsection*{Example 3: Given \( v = 2xy \)}

\textbf{Step 1 --- Verify harmonicity:}
\[
\frac{\partial^2 v}{\partial x^2} = 0, \qquad \frac{\partial^2 v}{\partial y^2} = 0, \qquad \nabla^2 v = 0. \quad \checkmark
\]
\textbf{Step 2 --- Compute partials (v-case):}
\[
\psi_1 = \frac{\partial v}{\partial y} = 2x, \qquad \psi_2 = \frac{\partial v}{\partial x} = 2y.
\]
\textbf{Step 3 --- Substitute \( x=z,\; y=0 \):}
\[
\psi_1(z,0) = 2z, \qquad \psi_2(z,0) = 0.
\]
\textbf{Step 4:}
\[
f'(z) = 2z + i(0) = 2z.
\]
\textbf{Step 5:}
\[
f(z) = \int 2z\,dz + C = z^2 + C.
\]
\textbf{Verification:} \( z^2 = x^2-y^2 + 2ixy \). Imaginary part \( = 2xy = v \). \checkmark

\begin{learnbox}{Result}
\( f(z) = z^2 + C \). Both \( u = x^2-y^2 \) and \( v=2xy \) lead to the same analytic function --- they are harmonic conjugates.
\end{learnbox}

% ---- Example 4 ----
\subsection*{Example 4: Rational Form --- Given \( u = \dfrac{x}{x^2+y^2} \)}

\textbf{Step 1 --- Verify harmonicity:}
\[
u = \frac{x}{x^2+y^2}.
\]
\[
\frac{\partial u}{\partial x} = \frac{(x^2+y^2) - x\cdot 2x}{(x^2+y^2)^2} = \frac{y^2-x^2}{(x^2+y^2)^2}.
\]
\[
\frac{\partial^2 u}{\partial x^2} = \frac{-2x(x^2+y^2)^2 - (y^2-x^2)\cdot 2(x^2+y^2)\cdot 2x}{(x^2+y^2)^4}
= \frac{2x(x^2-3y^2)}{(x^2+y^2)^3}.
\]
\[
\frac{\partial u}{\partial y} = \frac{-2xy}{(x^2+y^2)^2}.
\]
\[
\frac{\partial^2 u}{\partial y^2} = \frac{-2x(x^2+y^2)^2 + 2xy \cdot 2(x^2+y^2)\cdot 2y}{(x^2+y^2)^4}
= \frac{2x(3y^2-x^2)}{(x^2+y^2)^3}.
\]
\[
\nabla^2 u = \frac{2x(x^2-3y^2)}{(x^2+y^2)^3} + \frac{2x(3y^2-x^2)}{(x^2+y^2)^3} = 0. \quad \checkmark
\]
\textbf{Step 2:}
\[
\phi_1 = \frac{y^2-x^2}{(x^2+y^2)^2}, \qquad \phi_2 = \frac{-2xy}{(x^2+y^2)^2}.
\]
\textbf{Step 3 --- Substitute \( x=z,\; y=0 \):}
\[
\phi_1(z,0) = \frac{0 - z^2}{z^4} = \frac{-1}{z^2}, \qquad \phi_2(z,0) = 0.
\]
\textbf{Step 4:}
\[
f'(z) = -\frac{1}{z^2}.
\]
\textbf{Step 5:}
\[
f(z) = \int -\frac{1}{z^2}\,dz + C = \frac{1}{z} + C.
\]
\textbf{Verification:} \( \dfrac{1}{z} = \dfrac{\bar{z}}{|z|^2} = \dfrac{x-iy}{x^2+y^2} \). Real part \( = \dfrac{x}{x^2+y^2} = u \). \checkmark

\begin{learnbox}{Result}
\( f(z) = 1/z + C \). Rational harmonic functions often lead to rational analytic functions.
\end{learnbox}

% ---- Example 5 ----
\subsection*{Example 5: Given \( u = x^3 - 3xy^2 \)}

\textbf{Step 1:}
\[
\frac{\partial^2 u}{\partial x^2} = 6x, \qquad \frac{\partial^2 u}{\partial y^2} = -6x, \qquad \nabla^2 u = 0. \quad \checkmark
\]
\textbf{Step 2:}
\[
\phi_1 = 3x^2 - 3y^2, \qquad \phi_2 = -6xy.
\]
\textbf{Step 3:}
\[
\phi_1(z,0) = 3z^2, \qquad \phi_2(z,0) = 0.
\]
\textbf{Step 4:}
\[
f'(z) = 3z^2.
\]
\textbf{Step 5:}
\[
f(z) = z^3 + C.
\]
\textbf{Verification:} \( z^3 = (x+iy)^3 = x^3-3xy^2 + i(3x^2y-y^3) \). Real part \( = x^3-3xy^2 \). \checkmark

\begin{learnbox}{Result}
\( f(z) = z^3 + C \). Polynomial harmonics always yield the matching power of \( z \).
\end{learnbox}

% ---- Example 6 ----
\subsection*{Example 6: Given \( v = e^x \sin y \)}

\textbf{Step 1:}
\[
\frac{\partial^2 v}{\partial x^2} = e^x\sin y, \qquad \frac{\partial^2 v}{\partial y^2} = -e^x\sin y, \qquad \nabla^2 v = 0. \quad \checkmark
\]
\textbf{Step 2 (v-case):}
\[
\psi_1 = \frac{\partial v}{\partial y} = e^x\cos y, \qquad \psi_2 = \frac{\partial v}{\partial x} = e^x\sin y.
\]
\textbf{Step 3:}
\[
\psi_1(z,0) = e^z, \qquad \psi_2(z,0) = 0.
\]
\textbf{Step 4:}
\[
f'(z) = e^z + i(0) = e^z.
\]
\textbf{Step 5:}
\[
f(z) = e^z + C.
\]
\textbf{Verification:} Imaginary part of \( e^z \) is \( e^x \sin y = v \). \checkmark

\begin{learnbox}{Result}
\( f(z) = e^z + C \). Whether we start from \( u = e^x\cos y \) or \( v = e^x \sin y \), the analytic function is \( e^z \).
\end{learnbox}

% ---- Example 7 ----
\subsection*{Example 7: Given \( u = \ln\!\sqrt{x^2+y^2} \)}

Note \( u = \tfrac{1}{2}\ln(x^2+y^2) \).

\textbf{Step 1:}
\[
\frac{\partial^2 u}{\partial x^2} = \frac{y^2-x^2}{(x^2+y^2)^2}, \qquad \frac{\partial^2 u}{\partial y^2} = \frac{x^2-y^2}{(x^2+y^2)^2}, \qquad \nabla^2 u = 0. \quad \checkmark
\]
\textbf{Step 2:}
\[
\phi_1 = \frac{x}{x^2+y^2}, \qquad \phi_2 = \frac{y}{x^2+y^2}.
\]
\textbf{Step 3:}
\[
\phi_1(z,0) = \frac{z}{z^2} = \frac{1}{z}, \qquad \phi_2(z,0) = 0.
\]
\textbf{Step 4:}
\[
f'(z) = \frac{1}{z}.
\]
\textbf{Step 5:}
\[
f(z) = \ln z + C.
\]
\textbf{Verification:} \( \ln z = \ln|z| + i\arg(z) = \ln\sqrt{x^2+y^2} + i\arctan(y/x) \). Real part equals \( u \). \checkmark

\begin{learnbox}{Result}
\( f(z) = \ln z + C \). Logarithmic harmonics connect naturally to the complex logarithm.
\end{learnbox}

% ---- Example 8 (comparison) ----
\subsection*{Example 8: Given \( u = x^2 - y^2 + 2x \) --- Milne Thomson vs.\ Harmonic Conjugate Method}

\textbf{Milne Thomson approach:}

Step 1: \( u_{xx} = 2,\; u_{yy} = -2,\; \nabla^2 u = 0 \). \checkmark

Step 2: \( \phi_1 = 2x+2,\; \phi_2 = -2y \).

Step 3: \( \phi_1(z,0) = 2z+2,\; \phi_2(z,0) = 0 \).

Step 4: \( f'(z) = 2z + 2 \).

Step 5:
\[
f(z) = z^2 + 2z + C.
\]

\vspace{6pt}
\textbf{Harmonic Conjugate approach (for comparison):}

From C-R: \( v_y = u_x = 2x+2 \Rightarrow v = (2x+2)y + g(x) \).

Also \( v_x = -u_y = 2y \Rightarrow \dfrac{\partial}{\partial x}[(2x+2)y + g(x)] = 2y + g'(x) = 2y \Rightarrow g'(x) = 0 \Rightarrow g = c_1 \).

So \( v = 2xy + 2y + c_1 \), and:
\[
f(z) = (x^2-y^2+2x) + i(2xy+2y) + c_1 i = z^2 + 2z + C.
\]

Both methods give \( f(z) = z^2 + 2z + C \). \checkmark

\begin{learnbox}{Result}
\( f(z) = z^2 + 2z + C \). Milne Thomson required only a substitution and one integration; the harmonic conjugate method needed partial integration and a consistency check. Both are correct.
\end{learnbox}

% ------------------------------------------------------------------
% SECTION 7: COMPARISON (MANDATORY UPGRADE)
% ------------------------------------------------------------------
\section{Comparison with Harmonic Conjugate Method (MANDATORY UPGRADE)}

\begin{center}
\begin{tabular}{@{}p{3cm}p{4.5cm}p{4.5cm}p{3.5cm}@{}}
\toprule
\textbf{Method} & \textbf{Advantages} & \textbf{Limitations} & \textbf{Best Use Case} \\
\midrule
Harmonic Conjugate & Works from first principles; makes \( u \) and \( v \) explicit at every step & Requires integration of two expressions; constant-matching can be error-prone & When you need \( v \) explicitly (e.g., streamlines in fluid flow) \\
\midrule
Milne Thomson & No intermediate integration of \( u \) or \( v \) against \( y \) or \( x \); gives \( f(z) \) directly & Does not give \( v \) explicitly without expanding \( f(z) \); requires harmonicity check first & Exam problems asking for \( f(z) \) quickly; all standard textbook examples \\
\bottomrule
\end{tabular}
\end{center}

\begin{learnbox}{Choosing the Method}
Use \textbf{Milne Thomson} when the problem asks for \( f(z) \) in terms of \( z \). Use the \textbf{harmonic conjugate} method when the problem explicitly asks you to find the conjugate function \( v \) or when \( f(z) \) is not easily recognisable after integration.
\end{learnbox}

% ------------------------------------------------------------------
% SECTION 8: FLOWCHART (MANDATORY UPGRADE)
% ------------------------------------------------------------------
\section{Flowchart / Algorithm Summary (MANDATORY UPGRADE)}

\begin{center}
\begin{tikzpicture}[
  node distance=1.2cm,
  box/.style={rectangle, draw=blue!70!black, rounded corners, fill=blue!8, minimum width=7cm, minimum height=0.75cm, align=center, font=\small},
  dec/.style={diamond, draw=orange!80!black, fill=yellow!15, aspect=3, minimum width=6cm, align=center, font=\small},
  arr/.style={->, thick, draw=gray!60}
]
\node[box] (start) {START: Given real part \( u(x,y) \) or imaginary part \( v(x,y) \)};
\node[dec, below=of start] (check) {Compute \( \nabla^2(\text{given part}) \). Is it \( = 0 \)?};
\node[box, right=2cm of check] (stop) {STOP: Not harmonic. No analytic function exists.};
\node[dec, below=1.4cm of check] (which) {Which is given: \( u \) or \( v \)?};
\node[box, below left=1.2cm and 0.2cm of which] (ucalc) {Compute \( \phi_1 = u_x \),\; \( \phi_2 = u_y \)};
\node[box, below right=1.2cm and 0.2cm of which] (vcalc) {Compute \( \psi_1 = v_y \),\; \( \psi_2 = v_x \)};
\node[box, below=3.8cm of which] (sub) {Substitute \( x \to z,\; y \to 0 \)};
\node[box, below=1.2cm of sub] (form) {\( u \)-case: \( f'(z) = \phi_1(z,0) - i\phi_2(z,0) \)\quad\( v \)-case: \( f'(z) = \psi_1(z,0) + i\psi_2(z,0) \)};
\node[box, below=1.2cm of form] (integ) {Integrate \( f'(z) \) w.r.t.\ \( z \); add constant \( C \)};
\node[box, below=1.2cm of integ] (verify) {Verify: expand \( f(z) \) and check real/imaginary part matches};
\node[box, below=1.2cm of verify] (end) {DONE: \( f(z) \) is the required analytic function};

\draw[arr] (start) -- (check);
\draw[arr] (check) -- node[left, font=\small]{Yes} (which);
\draw[arr] (check) -- node[above, font=\small]{No} (stop);
\draw[arr] (which) -- node[left, font=\small]{\( u \)} (ucalc);
\draw[arr] (which) -- node[right, font=\small]{\( v \)} (vcalc);
\draw[arr] (ucalc) |- (sub);
\draw[arr] (vcalc) |- (sub);
\draw[arr] (sub) -- (form);
\draw[arr] (form) -- (integ);
\draw[arr] (integ) -- (verify);
\draw[arr] (verify) -- (end);
\end{tikzpicture}
\end{center}

\begin{learnbox}{Using the Flowchart}
The two decision diamonds are the critical checkpoints: (1) harmonicity must hold before anything, and (2) the sign pattern in \( f'(z) \) depends on whether \( u \) or \( v \) is given.
\end{learnbox}

% ------------------------------------------------------------------
% SECTION 9: EXCEL EXAMPLE (MANDATORY)
% ------------------------------------------------------------------
\section{Excel Example (MANDATORY)}

\textbf{Problem:} Given \( u(x,y) = x^2 - y^2 \), use a spreadsheet to compute \( \phi_1 = u_x = 2x \), \( \phi_2 = u_y = -2y \), apply the substitution \( y = 0, x = z \) (treating \( z \) as real here), and observe how \( f'(z) = 2z \).

\begin{center}
\begin{tabular}{@{}ccccccc@{}}
\toprule
\textbf{Col A} & \textbf{Col B} & \textbf{Col C} & \textbf{Col D} & \textbf{Col E} & \textbf{Col F} & \textbf{Col G} \\
\( x \) & \( y \) & \( u = x^2-y^2 \) & \( \phi_1 = u_x \) & \( \phi_2 = u_y \) & \( \phi_1(z,0) \) & \( f'(z) = \phi_1(z,0) \) \\
\midrule
1 & 0 & 1 & 2 & 0 & 2 & 2 \\
2 & 0 & 4 & 4 & 0 & 4 & 4 \\
3 & 0 & 9 & 6 & 0 & 6 & 6 \\
4 & 0 & 16 & 8 & 0 & 8 & 8 \\
5 & 0 & 25 & 10 & 0 & 10 & 10 \\
\midrule
1 & 1 & 0 & 2 & $-2$ & --- & --- \\
2 & 1 & 3 & 4 & $-2$ & --- & --- \\
\bottomrule
\end{tabular}
\end{center}

\textbf{Excel Formulas used (row 2 as example):}
\begin{itemize}
  \item \texttt{C2} = \texttt{=A2\^{}2 - B2\^{}2}
  \item \texttt{D2} = \texttt{=2*A2} \quad (this is \( u_x = 2x \))
  \item \texttt{E2} = \texttt{=-2*B2} \quad (this is \( u_y = -2y \))
  \item \texttt{F2} = \texttt{=2*A2} \quad (substitute \( y=0 \), so \( \phi_1(z,0) = 2x \) with \( x = z \))
  \item \texttt{G2} = \texttt{=F2 - 0*E2} \quad (\( f'(z) = \phi_1 - i\phi_2 \); imaginary part is 0 here)
\end{itemize}

The last column confirms that \( f'(z) = 2z \) for \( z = 1,2,3,4,5 \). Integrating gives \( f(z) = z^2 + C \), consistent with the worked example.

\begin{learnbox}{Spreadsheet Strategy}
Columns D and E automate the derivative computation. Column F fixes \( y = 0 \) to complete the Milne Thomson substitution. This structure works for any \( u(x,y) \) by changing the formula in column C.
\end{learnbox}

% ------------------------------------------------------------------
% SECTION 10: PYTHON EXAMPLE (MANDATORY)
% ------------------------------------------------------------------
\section{Python Example (MANDATORY)}

The following complete Python script uses SymPy to automate the Milne Thomson procedure for \( u = e^x \cos y \).

\begin{lstlisting}[language=Python]
# Milne Thomson Method using SymPy
# Example: u = exp(x)*cos(y)

import sympy as sp

# Define symbols
x, y, z = sp.symbols('x y z')

# Define u(x, y)
u = sp.exp(x) * sp.cos(y)

# Step 1: Verify Laplace's equation
u_xx = sp.diff(u, x, 2)
u_yy = sp.diff(u, y, 2)
laplacian = sp.simplify(u_xx + u_yy)
print("Laplacian of u:", laplacian)
# Expected output: Laplacian of u: 0

# Step 2: Compute phi1 and phi2
phi1 = sp.diff(u, x)   # partial u / partial x
phi2 = sp.diff(u, y)   # partial u / partial y
print("phi1 =", phi1)  # Expected: exp(x)*cos(y)
print("phi2 =", phi2)  # Expected: -exp(x)*sin(y)

# Step 3: Substitute x -> z, y -> 0
phi1_sub = phi1.subs([(x, z), (y, 0)])
phi2_sub = phi2.subs([(x, z), (y, 0)])
print("phi1(z,0) =", phi1_sub)  # Expected: exp(z)
print("phi2(z,0) =", phi2_sub)  # Expected: 0

# Step 4: Form f'(z)
f_prime = phi1_sub - sp.I * phi2_sub
print("f'(z) =", f_prime)  # Expected: exp(z)

# Step 5: Integrate to get f(z)
f_z = sp.integrate(f_prime, z)
print("f(z) =", f_z, "+ C")   # Expected: f(z) = exp(z) + C

# Step 6: Verification - expand f(z) and extract real part
# Rewrite in terms of x, y by substituting z = x + I*y
z_val = x + sp.I * y
f_expanded = sp.exp(z_val)
real_part = sp.re(f_expanded)
real_part_simplified = sp.simplify(real_part)
print("Real part of f(z):", real_part_simplified)
# Expected: Real part of f(z): exp(x)*cos(y)
# This matches u(x, y) = exp(x)*cos(y). Verification complete.
\end{lstlisting}

\begin{learnbox}{Python Takeaway}
SymPy's \texttt{diff}, \texttt{subs}, and \texttt{integrate} mirror the four steps of the Milne Thomson algorithm exactly. The substitution step \texttt{subs([(x,z),(y,0)])} is the computational heart of the method.
\end{learnbox}

% ------------------------------------------------------------------
% SECTION 11: VIVA QUESTIONS
% ------------------------------------------------------------------
\section{Viva Questions}

\begin{enumerate}
  \item \textbf{What is the Milne Thomson method and what problem does it solve?}\\
  The Milne Thomson method constructs an analytic function \( f(z) \) directly from a given harmonic function \( u(x,y) \) or \( v(x,y) \) by forming \( f'(z) \) using the C-R equations and then applying the substitution \( x = z,\; y = 0 \) followed by integration.

  \item \textbf{Why must we check that \( u \) is harmonic before applying the Milne Thomson method?}\\
  An analytic function exists only if its real and imaginary parts satisfy Laplace's equation. If \( u \) is not harmonic (\( \nabla^2 u \neq 0 \)), no analytic function can have \( u \) as its real part, and the method would yield a meaningless result.

  \item \textbf{Explain why substituting \( x = z,\; y = 0 \) is valid.}\\
  The formula \( f'(z) = \phi_1(x,y) - i\phi_2(x,y) \) holds for all \( z = x+iy \). Along the real axis \( y=0 \), the complex variable equals the real variable \( x \). By analytic continuation, the expression \( \phi_1(z,0) - i\phi_2(z,0) \) equals \( f'(z) \) throughout the domain of analyticity.

  \item \textbf{Write the formula for \( f'(z) \) when \( v(x,y) \) is given.}\\
  \( f'(z) = \dfrac{\partial v}{\partial y}\bigg|_{(z,0)} + i\,\dfrac{\partial v}{\partial x}\bigg|_{(z,0)} \).

  \item \textbf{What is the role of the arbitrary constant \( C \) in the Milne Thomson method?}\\
  The constant \( C \) represents an arbitrary complex constant arising from the indefinite integration of \( f'(z) \). It accounts for all analytic functions sharing the same real (or imaginary) part; typically \( C \) is fixed by an initial condition such as \( f(0) = 0 \).

  \item \textbf{Given \( u = \ln\sqrt{x^2+y^2} \), what is \( f(z) \)?}\\
  \( f(z) = \ln z + C \).

  \item \textbf{Can the Milne Thomson method be applied if the given function is not differentiable at isolated points?}\\
  The method applies in domains where \( u \) (or \( v \)) is harmonic and differentiable. Isolated singularities (where the function is undefined) must be excluded from the domain; the method still works in any simply connected domain that avoids those singularities.

  \item \textbf{If \( f(z) = z^3 + C \), what are the harmonic functions \( u \) and \( v \)?}\\
  Expanding: \( z^3 = (x+iy)^3 = x^3 - 3xy^2 + i(3x^2y - y^3) \). So \( u = x^3 - 3xy^2 \) and \( v = 3x^2y - y^3 \).
\end{enumerate}

% ------------------------------------------------------------------
% SECTION 12: DESCRIPTIVE QUESTIONS
% ------------------------------------------------------------------
\section{Descriptive Questions}

\begin{enumerate}
  \item State and derive the Milne Thomson formula for constructing an analytic function when the real part \( u(x,y) \) is given. Clearly justify the substitution \( x = z,\; y = 0 \). Apply the formula to find \( f(z) \) given \( u = x^3 - 3xy^2 \).

  \item Derive the Milne Thomson formula for the case when the imaginary part \( v(x,y) \) is given. Using this formula, construct \( f(z) \) when \( v = e^x \sin y \), and verify your answer by expanding \( f(z) \) in terms of \( x \) and \( y \).

  \item Compare the Milne Thomson method and the harmonic conjugate method for constructing analytic functions. Illustrate by solving the same problem \( u = x^2 - y^2 + 2x \) using both methods and confirm that both give the same analytic function.

  \item Given \( u = \dfrac{x}{x^2+y^2} \), verify that \( u \) is harmonic, then use the Milne Thomson method to find \( f(z) \). State any restrictions on the domain.

  \item A student applies the Milne Thomson method to \( u = x^2 + y^2 \) and gets \( f(z) = z^2 + C \). Identify the error and explain what goes wrong. What is \( \nabla^2(x^2+y^2) \)?
\end{enumerate}

% ------------------------------------------------------------------
% SECTION 13: PRACTICE PROBLEMS
% ------------------------------------------------------------------
\section{Practice Problems}

\begin{enumerate}
  \item Find \( f(z) \) given \( u = x^4 - 6x^2y^2 + y^4 \).\\
  \textit{Hint:} Compute \( u_x = 4x^3 - 12xy^2 \); after substitution, \( f'(z) = 4z^3 \); hence \( f(z) = z^4 + C \).

  \item Find \( f(z) \) given \( v = x^3y - xy^3 \).\\
  \textit{Hint:} \( v \)-case; \( \psi_1(z,0) = z^3 \), \( \psi_2(z,0) = 0 \); integrate to get \( f(z) = \frac{z^4}{4} \cdot i + C \). [Check sign carefully.]

  \item Find \( f(z) \) given \( u = e^{2x}\cos 2y \).\\
  \textit{Hint:} \( \phi_1(z,0) = 2e^{2z} \); \( f(z) = e^{2z} + C \).

  \item Find \( f(z) \) given \( v = \arctan(y/x) \) (principal value, \( x > 0 \)).\\
  \textit{Hint:} \( \psi_1 = v_y = x/(x^2+y^2) \); \( \psi_1(z,0) = 1/z \); \( f(z) = \ln z + C \).

  \item Verify that \( u = \sin x \cosh y \) is harmonic and find the analytic function \( f(z) \).\\
  \textit{Hint:} \( \phi_1(z,0) = \cos z \); \( f(z) = \sin z + C \).

  \item Given \( u = 2x(1-y) \), verify harmonicity and find \( f(z) \).\\
  \textit{Hint:} \( u_{xx} = 0,\; u_{yy} = 0 \); \( \phi_1 = 2(1-y),\; \phi_1(z,0) = 2 \); \( \phi_2 = -2x,\; \phi_2(z,0) = -2z \); \( f'(z) = 2 + 2iz \); \( f(z) = 2z + iz^2 + C \).
\end{enumerate}

% ------------------------------------------------------------------
% SECTION 14: MCQs
% ------------------------------------------------------------------
\section{MCQs}

\begin{enumerate}
  \item In the Milne Thomson method, the substitution made in \( f'(z) = \phi_1(x,y) - i\phi_2(x,y) \) is:
  \begin{enumerate}[(a)]
    \item \( x = 0,\; y = z \)
    \item \textbf{(b)} \( x = z,\; y = 0 \)
    \item \( x = iz,\; y = z \)
    \item \( x = 1/z,\; y = 1 \)
  \end{enumerate}
  \textit{Explanation:} The real axis corresponds to \( y = 0 \), making \( z = x \); analytic continuation extends the result to the whole domain.

  \item If \( u = x^2 - y^2 \), the analytic function \( f(z) \) by Milne Thomson is:
  \begin{enumerate}[(a)]
    \item \( z + C \)
    \item \( 2z + C \)
    \item \textbf{(c)} \( z^2 + C \)
    \item \( z^3 + C \)
  \end{enumerate}
  \textit{Explanation:} \( \phi_1 = 2x \to 2z \); integrate \( 2z \) to get \( z^2 + C \).

  \item The formula for \( f'(z) \) when \( v(x,y) \) is given is:
  \begin{enumerate}[(a)]
    \item \( v_x - i\,v_y \)
    \item \( v_y - i\,v_x \)
    \item \textbf{(c)} \( v_y + i\,v_x \)
    \item \( v_x + i\,v_y \)
  \end{enumerate}
  \textit{Explanation:} Using C-R equations: \( f'(z) = u_x + iv_x = v_y + iv_x \); the \( v \)-given formula uses \( +i \).

  \item Before applying Milne Thomson, one must verify that the given function:
  \begin{enumerate}[(a)]
    \item Is continuous everywhere
    \item \textbf{(b)} Satisfies Laplace's equation
    \item Is bounded in the domain
    \item Has continuous partial derivatives up to 4th order
  \end{enumerate}
  \textit{Explanation:} The real (or imaginary) part of an analytic function must be harmonic, i.e., it must satisfy \( \nabla^2 \phi = 0 \).

  \item Given \( u = e^x \cos y \), the analytic function \( f(z) \) is:
  \begin{enumerate}[(a)]
    \item \( e^{iz} + C \)
    \item \( \cos z + C \)
    \item \( ie^z + C \)
    \item \textbf{(d)} \( e^z + C \)
  \end{enumerate}
  \textit{Explanation:} \( \phi_1(z,0) = e^z,\; \phi_2(z,0) = 0 \); integrating \( e^z \) gives \( e^z + C \).
\end{enumerate}

% ------------------------------------------------------------------
% SECTION 15: COMMON MISTAKES
% ------------------------------------------------------------------
\section{Common Mistakes}

\begin{mistakebox}{Common Mistakes in Milne Thomson Method}
\begin{tabular}{@{}p{4.5cm}p{5.5cm}p{5cm}@{}}
\toprule
\textbf{Mistake} & \textbf{What Goes Wrong} & \textbf{Correct Approach} \\
\midrule
Wrong substitution: \( y=z,\; x=0 \) instead of \( x=z,\; y=0 \) & The formula derives \( f'(z) \) along the real axis where \( y=0 \) and \( x \) plays the role of \( z \); reversing them gives a nonsensical function & Always substitute \( x \to z \) and \( y \to 0 \) simultaneously \\
\midrule
Sign error in \( v \)-given formula & Writing \( f'(z) = \psi_1 - i\psi_2 \) instead of \( \psi_1 + i\psi_2 \) & In the \( v \)-case, use \( f'(z) = v_y + iv_x \) (note \( +i \)) \\
\midrule
Forgetting the integration constant \( C \) & The solution is incomplete; boundary or initial conditions cannot be applied & Always write \( f(z) = \ldots + C \) after integrating \\
\midrule
Applying the method without checking harmonicity & If \( \nabla^2 u \neq 0 \), no analytic \( f \) with that real part exists; the result will be wrong & Compute \( u_{xx} + u_{yy} \) first; proceed only if it equals 0 \\
\midrule
Not verifying the answer & A computational slip (e.g., dropping a factor) may go undetected & After finding \( f(z) \), expand and confirm that real or imaginary part matches the given function \\
\midrule
Confusing \( \phi_1 \) with \( \phi_2 \) & Swapping \( u_x \) and \( u_y \) in the formula changes the sign of the imaginary part and gives the wrong \( f(z) \) & Define \( \phi_1 = u_x \) and \( \phi_2 = u_y \) explicitly before substituting \\
\bottomrule
\end{tabular}
\end{mistakebox}

% ------------------------------------------------------------------
% SECTION 16: QUICK RECAP
% ------------------------------------------------------------------
\section{Quick Recap}

\begin{learnbox}{Quick Recap --- Milne Thomson Method}
\begin{itemize}
  \item \textbf{Core idea:} Given \( u \) or \( v \), construct \( f'(z) \) using C-R equations, substitute \( x=z,\;y=0 \), integrate to get \( f(z) \).
  \item \textbf{u-given formula:} \( f'(z) = \phi_1(z,0) - i\,\phi_2(z,0) \) where \( \phi_1 = u_x,\;\phi_2 = u_y \).
  \item \textbf{v-given formula:} \( f'(z) = \psi_1(z,0) + i\,\psi_2(z,0) \) where \( \psi_1 = v_y,\;\psi_2 = v_x \).
  \item \textbf{Mandatory first step:} Always verify \( \nabla^2(\text{given part}) = 0 \) before anything else.
  \item \textbf{Key substitution:} Setting \( y=0,\;x=z \) converts a two-variable expression into a function of \( z \) alone, enabling direct integration.
  \item \textbf{Advantage over harmonic conjugate method:} No intermediate integration against \( y \) or \( x \); \( f(z) \) is obtained in one integration step.
  \item \textbf{Always add:} An arbitrary complex constant \( C \) after integration.
  \item \textbf{Standard results:} \( u = x^n \)-type \( \to f(z)=z^n+C \); \( u=e^x\cos y \to f(z)=e^z+C \); \( u=\ln\sqrt{x^2+y^2} \to f(z)=\ln z + C \).
\end{itemize}
\end{learnbox}

\end{document}
